\chapter{Machine Learning}

\section{Supervised Learning}
\dfn{Supervised Learning}{
Learn a mapping $x \rightarrow y$ from labeled data.
}
\dfn{Classification}{
Predict discrete labels (e.g., spam/ham). Metrics: accuracy, precision/recall/F1, ROC-AUC.
}
\dfn{Regression}{
Predict continuous targets. Metrics: MAE, MSE/RMSE, $R^2$.
}
\dfn{Bias--Variance}{
High bias underfits; high variance overfits. Use regularization, more data, or simpler/ensemble models.
}
\subsection*{Common Models}
\begin{itemize}
    \item Logistic regression (classification), linear regression (OLS/Ridge/Lasso).
    \item $k$-NN, Naive Bayes, Decision Trees, Random Forests, Gradient Boosting (XGBoost/GBM).
    \item SVM (linear/RBF kernels).
\end{itemize}

\subsection*{Sklearn Pipeline (train/test split)}
\begin{lstlisting}
from sklearn.model_selection import train_test_split
from sklearn.preprocessing import StandardScaler
from sklearn.pipeline import make_pipeline
from sklearn.linear_model import LogisticRegression
from sklearn.metrics import accuracy_score, classification_report

X_train, X_test, y_train, y_test = train_test_split(X, y, test_size=0.2, random_state=42, stratify=y)

clf = make_pipeline(StandardScaler(), LogisticRegression(max_iter=1000))
clf.fit(X_train, y_train)
pred = clf.predict(X_test)
print("acc", accuracy_score(y_test, pred))
print(classification_report(y_test, pred))
\end{lstlisting}

\subsection*{Cross-Validation \& Grid Search}
\begin{lstlisting}
from sklearn.model_selection import GridSearchCV
params = {"logisticregression__C":[0.1,1,10]}
grid = GridSearchCV(clf, params, cv=5, scoring="f1")
grid.fit(X, y)
print(grid.best_params_, grid.best_score_)
\end{lstlisting}

\section{Unsupervised Learning}
\dfn{Unsupervised Learning}{
Discover structure in unlabeled data.
}
\dfn{Clustering}{
Group similar instances (e.g., K-Means, Hierarchical, DBSCAN).
}
\dfn{Dimensionality Reduction}{
Project data to lower dimensions while preserving variance/structure (PCA, t-SNE/UMAP for viz).
}
\dfn{Silhouette Score}{
Internal cluster quality: cohesion vs separation; range $[-1,1]$.
}

\subsection*{K-Means Example}
\begin{lstlisting}
from sklearn.cluster import KMeans
from sklearn.metrics import silhouette_score

kmeans = KMeans(n_clusters=3, random_state=0, n_init="auto")
labels = kmeans.fit_predict(X)
print("silhouette", silhouette_score(X, labels))
\end{lstlisting}

\subsection*{PCA Example}
\begin{lstlisting}
from sklearn.decomposition import PCA
pca = PCA(n_components=2, whiten=True, random_state=0)
X_2d = pca.fit_transform(X)
print("explained var", pca.explained_variance_ratio_)
\end{lstlisting}

\section{Other ML Libraries}
\dfn{scikit-learn}{
General-purpose Python ML: preprocessing, model selection, classic algorithms.
}
\dfn{TensorFlow / Keras}{
Deep learning framework + high-level API for building neural nets.
}
\dfn{PyTorch}{
Dynamic graph deep learning library, strong research adoption.
}
\dfn{XGBoost / LightGBM}{
Optimized gradient boosting for tabular data.
}

\subsection*{Train a Simple MLP (Keras)}
\begin{lstlisting}
from tensorflow import keras
model = keras.Sequential([
    keras.layers.Input(shape=(X_train.shape[1],)),
    keras.layers.Dense(64, activation="relu"),
    keras.layers.Dense(1, activation="sigmoid")
])
model.compile(optimizer="adam", loss="binary_crossentropy", metrics=["accuracy"])
model.fit(X_train, y_train, epochs=10, batch_size=32, validation_split=0.1)
\end{lstlisting}

\section{Neural Networks (Detailed)}
\dfn{Perceptron / Neuron}{
Linear transform plus nonlinearity: $h = \sigma(Wx + b)$; nonlinearity enables universal approximation.
}
\dfn{Activation Functions}{
ReLU ($\max(0,x)$), Leaky ReLU, GELU, Sigmoid (bounded, can saturate), Tanh (zero-centered).
}
\dfn{Loss Functions}{
Classification: cross-entropy; regression: MSE/MAE. Optimization minimizes empirical risk.
}
\dfn{Backpropagation}{
Chain rule to compute gradients of loss w.r.t. parameters layer by layer.
}
\dfn{Gradient-Based Optimizers}{
SGD (with/without momentum), RMSProp, Adam (adaptive moments).
}
\dfn{Regularization}{
Weight decay ($L_2$), dropout, early stopping, data augmentation, batch/ layer normalization.
}
\dfn{Initialization}{
Xavier/Glorot for tanh; He for ReLU family to keep variance stable across layers.
}
\dfn{Batching}{
Mini-batch training stabilizes gradients; batch size trades compute vs noise.
}
\dfn{Vanishing / Exploding Gradients}{
Deep nets can suffer when derivatives shrink or blow up; mitigations: better inits, ReLU-like activations, normalization, residual connections.
}

\subsection*{Feedforward MLP (PyTorch)}
\begin{lstlisting}
import torch
from torch import nn

class MLP(nn.Module):
    def __init__(self, d_in, d_hidden=128):
        super().__init__()
        self.net = nn.Sequential(
            nn.Linear(d_in, d_hidden),
            nn.ReLU(),
            nn.Dropout(0.2),
            nn.Linear(d_hidden, 1),
            nn.Sigmoid()
        )
    def forward(self, x):
        return self.net(x)

model = MLP(d_in=X_train.shape[1])
loss_fn = nn.BCELoss()
opt = torch.optim.Adam(model.parameters(), lr=1e-3)

for epoch in range(10):
    model.train()
    xb = torch.tensor(X_train, dtype=torch.float32)
    yb = torch.tensor(y_train, dtype=torch.float32).unsqueeze(1)
    pred = model(xb)
    loss = loss_fn(pred, yb)
    opt.zero_grad()
    loss.backward()
    opt.step()
    print(epoch, float(loss))
\end{lstlisting}

\subsection*{Feedforward MLP (Keras)}
\begin{lstlisting}
from tensorflow import keras
model = keras.Sequential([
    keras.layers.Input(shape=(X_train.shape[1],)),
    keras.layers.Dense(256, activation="relu", kernel_initializer="he_normal"),
    keras.layers.Dropout(0.2),
    keras.layers.Dense(128, activation="relu"),
    keras.layers.Dense(1, activation="sigmoid")
])
model.compile(optimizer=keras.optimizers.Adam(1e-3),
              loss="binary_crossentropy",
              metrics=["accuracy"])
model.fit(X_train, y_train, epochs=15, batch_size=32,
          validation_split=0.2, callbacks=[
              keras.callbacks.EarlyStopping(patience=3, restore_best_weights=True)
          ])
\end{lstlisting}

\subsection*{Convolutional Neural Networks (CNN) Essentials}
\dfn{Convolution}{
Local receptive field with shared weights; parameters: kernel size, stride, padding, channels.
}
\dfn{Pooling}{
Downsample (max/avg) to reduce spatial dims and add invariance.
}
\dfn{Padding / Stride}{
Padding preserves spatial size; stride controls step size (stride>1 downsamples).
}

\subsection*{Simple CNN (Keras)}
\begin{lstlisting}
model = keras.Sequential([
    keras.layers.Input(shape=(64,64,3)),
    keras.layers.Conv2D(32, 3, padding="same", activation="relu"),
    keras.layers.MaxPool2D(),
    keras.layers.Conv2D(64, 3, padding="same", activation="relu"),
    keras.layers.MaxPool2D(),
    keras.layers.Flatten(),
    keras.layers.Dense(128, activation="relu"),
    keras.layers.Dense(10, activation="softmax")
])
model.compile(optimizer="adam",
              loss="sparse_categorical_crossentropy",
              metrics=["accuracy"])
\end{lstlisting}

\subsection*{Recurrent / Sequence Models}
\dfn{RNN / LSTM / GRU}{
Handle sequences by maintaining hidden state; LSTM/GRU mitigate vanishing gradients with gates.
}

\begin{lstlisting}
model = keras.Sequential([
    keras.layers.Input(shape=(None, d_embed)),
    keras.layers.GRU(128, return_sequences=False),
    keras.layers.Dense(1, activation="sigmoid")
])
model.compile(optimizer="adam", loss="binary_crossentropy", metrics=["accuracy"])
\end{lstlisting}

\subsection*{Transformers (concept)}
\dfn{Self-Attention}{
Computes token-to-token weights to mix information without recurrence or convolution.
}
\dfn{Positional Encoding}{
Adds order information for sequences in transformer models.
}

\section{Practice / Exam Prompts (NN-Focused)}
\pr{Init + activation pairing}{
Why pair He init with ReLU? To preserve variance through layers and reduce vanishing gradients.
}
\pr{Regularization plan}{
Given overfitting CNN: add dropout, data augmentation, weight decay, and early stopping; maybe reduce depth.
}
\pr{Optimizer choice}{
When would you prefer SGD with momentum over Adam? For better generalization on vision tasks; Adam for faster early convergence or sparse grads.
}
\pr{Batch norm effect}{
Explain how batch norm stabilizes training (normalizes intermediate activations, allows higher learning rates).
}
\pr{Sequence model pick}{
Short text classification with little data: use GRU; long contexts or parallel compute: prefer transformer encoder.
}

\section{Practice / Exam Prompts}
\pr{Model selection}{
Given a tabular binary classification task, outline: split data, scale features, baseline logistic regression, try tree/ensemble, compare ROC-AUC, tune with CV.
}
\pr{Overfitting diagnosis}{
Train/val gap high $\rightarrow$ high variance. Mitigate with regularization, fewer parameters, more data, early stopping, dropout (DL).
}
\pr{Clustering choice}{
When clusters are spherical and k is known, use K-Means; if density-based and arbitrary shapes, use DBSCAN; evaluate with silhouette.
}
\pr{PCA intuition}{
PCA finds orthogonal directions of max variance; good for noise reduction and visualization; watch scaling before PCA.
}
\pr{Boosting vs bagging}{
Bagging (RF) reduces variance via parallel trees; boosting (GBM/XGB) reduces bias via sequential trees; boosting can overfit noise if too deep/too many rounds.
}
